\documentclass[12pt,a4paper]{article}
\usepackage[utf8]{inputenc}
\usepackage[T1]{fontenc}
\usepackage{lmodern}
\usepackage[ngerman]{babel}
\usepackage{amsmath,amssymb,amsthm}
\usepackage{graphicx}
\usepackage{hyperref}
\usepackage{url}
\usepackage{xcolor}
\usepackage{booktabs}
\usepackage{geometry}
\geometry{margin=2.5cm}

% Eigene Befehle
\newcommand{\Keff}{K_{\mathrm{eff}}}
\newcommand{\kappac}{\kappa_c}
\newcommand{\eps}{\varepsilon}
\newcommand{\betagamma}{\beta\gamma}
\newcommand{\betac}{\beta}
\newcommand{\tauf}{\tau}

\newtheorem{theorem}{Theorem}
\newtheorem{lemma}{Lemma}
\newtheorem{corollary}{Korollar}
\newtheorem{definition}{Definition}
\newtheorem{prediction}{Vorhersage}

\title{\textbf{Anhang S. Negative Zeitverzögerung von Photonen in kalten Atomen als Manifestation koordinierter Quantendynamik: \\ Eine überarbeitete Analyse im Rahmen der Yakushev Unified Coordination Theory (YUCT)}}
\author{Alexey V. Yakushev\textsuperscript{1} \\
	\textsuperscript{1}Yakushev Research, \\ 
	\textbf YUCT-Theorie: \url{https://yuct.org/}, \\ 
	\textbf YPSDC-Protokoll: \url{https://ypsdc.com/}\\
	\textbf DOI: \url{https://doi.org/10.5281/zenodo.18444598}\\
}
\date{Mai 2026 \\ \large V.3.0. \\(Korrigierte logarithmische Skalierung und Rücknahme der 13\%-Temperaturvorhersage)}

\begin{document}
	
	\maketitle
	\newpage
	\begin{abstract}
		Neueste bahnbrechende Experimente (Angulo et al., 2024) haben gezeigt, dass Photonen, die eine ultra-kalte Atomwolke durchqueren, das Medium mit einer negativen Zeitverzögerung verlassen können – der Peak des Wellenpakets erscheint früher, als es eingetreten ist. Dieses kontraintuitive Phänomen, das die Kausalität nicht verletzt, offenbart tiefgreifende Quanteninterferenzeffekte. In dieser Arbeit überdenken wir die Erklärung im Rahmen der Yakushev Unified Coordination Theory (YUCT). Aufbauend auf dem universellen Skalierungsgesetz \(\eps = \kappac \alpha (\ln\Keff)^{\betac}\) mit \(\betac = 2/3\) und \(\kappac = 1/3\) (abgeleitet in Anhang X) zeigen wir, dass die korrekte Beziehung zwischen der Gruppenverzögerung \(\tau_g\) und der Koordinationseffizienz \(\Keff\) logarithmisch ist: \(|\tau_g| = \alpha_0 \ln\Keff\). Die Existenz negativer Zeitverzögerungen wird bereits durch das Experiment bestätigt; Werte von \((-1,14 \pm 0,22)\tau_0\) wurden beobachtet. Im Gegensatz zu einem früheren linearen Ansatz impliziert die logarithmische Form jedoch, dass die Temperaturabhängigkeit von \(\tau_g\) extrem schwach ist. Unter Verwendung der kanonischen Temperaturdegradation \(\Keff(T) = K_0 (T/T_0)^{-\gamma}\) mit \(\gamma \approx 0,30\) (abgeleitet aus geophysikalischen Daten) finden wir, dass die relative Änderung von \(|\tau_g|\) bei Verdopplung der Temperatur in der Größenordnung von \(10^{-8}\) liegt, weit unterhalb der Empfindlichkeit aktueller Experimente. Folglich ist der zuvor behauptete 13\%-Effekt ein Artefakt eines falschen linearen Modells und muss zurückgenommen werden. Wir liefern den korrigierten theoretischen Rahmen und diskutieren alternative experimentelle Tests, die für die logarithmische Skalierung empfindlich sein könnten, wie etwa die Variation der optischen Tiefe oder der Pulsdauer. Die Kernvorhersage der YUCT – dass negative Gruppenverzögerungen eine Manifestation der Koordinationseffizienz sind – bleibt gültig, aber ihre Testbarkeit in kalten Atomsystemen erfordert verfeinerte Aufbauten.
	\end{abstract}
	
	\noindent\textbf{Schlüsselwörter:} YUCT, Koordinationseffizienz, negative Gruppenverzögerung, Photon-Atom-Wechselwirkung, logarithmische Skalierung, Temperaturunabhängigkeit, Rücknahme.
	
	\newpage
	\tableofcontents
	\newpage
	
	\section{Einleitung}
	
	Die Ausbreitung von Licht durch resonante Medien ist seit Jahrzehnten ein Eckpfeiler der Quantenoptik. Gruppenverzögerungen, Pulsumformung und überlichtschnelle Effekte sind theoretisch und experimentell gut verstanden \cite{Brillouin1960, Milonni2004}. Dennoch besteht eine grundlegende Frage: Was ist die physikalische Bedeutung der Zeit, die ein Photon als atomare Anregung verbringt?
	
	Ein bahnbrechendes Experiment von Angulo et al. \cite{Angulo2024} ging dieser Frage nach, indem es die durch ein transmittiertes Photon auf einen separaten Sondenstrahl induzierte Kreuz-Kerr-Phasenverschiebung maß. Sie fanden heraus, dass die integrierte atomare Anregungszeit \(\tau_T\) für ein transmittiertes Photon negativ sein kann und mit der Gruppenverzögerung \(\tau_g\) übereinstimmt, selbst wenn \(\tau_g < 0\) ist. Konkret berichtete das Experiment Werte von \(\tau_T/\tau_0 = -1,14 \pm 0,22\) (für einen 10-ns-Puls, OD~4) und \(0,54 \pm 0,28\) (für andere Parameter) \cite{Angulo2024, Nixon2024}. Dieses bemerkenswerte Ergebnis fordert einfache Intuitionen heraus und verlangt nach einem tieferen theoretischen Rahmen.
	
	Die Yakushev Unified Coordination Theory (YUCT) \cite{Yakushev2026a} bietet einen solchen Rahmen. YUCT postuliert, dass alle physikalischen Wechselwirkungen, von der Quantenoptik bis zur Gravitation, Manifestationen eines fundamentalen Koordinationsprozesses sind, der durch die \textbf{Koordinationseffizienz} \(\Keff\) quantifiziert wird. Ein universelles Skalierungsgesetz verknüpft jeden relativen Fehler \(\eps\) mit \(\Keff\) (siehe Anhang X, Gl.~\ref{eq:error-law-corrected}):
	\begin{equation}
		\eps = \kappac \alpha (\ln\Keff)^{\betac}, \qquad \betac = \frac{2}{3},\ \kappac = \frac{1}{3},
		\label{eq:scaling}
	\end{equation}
	wobei \(\alpha\) systemabhängig ist. Dieses Gesetz wurde über 40 Größenordnungen hinweg verifiziert, von DNA-Replikationsfehlern bis zu Fluktuationen des kosmischen Mikrowellenhintergrunds \cite{Yakushev2026c}. Im Kontext der Licht-Materie-Wechselwirkung misst \(\Keff\) die Kohärenz zwischen dem Photon und dem atomaren Ensemble.
	
	Hier zeigen wir, dass das Experiment zur negativen Zeitverzögerung eine direkte Konsequenz der YUCT ist. Allerdings müssen wir einen früheren Versuch korrigieren, der eine lineare Beziehung zwischen \(\tau_g\) und \(\Keff\) annahm. Die korrekte Beziehung, abgeleitet aus der informationstheoretischen Struktur der YUCT, ist logarithmisch. Dies ändert drastisch die vorhergesagte Temperaturabhängigkeit, die mit der aktuellen experimentellen Präzision vernachlässigbar wird. Wir ziehen daher die zuvor veröffentlichte 13\%-Vorhersage zurück und liefern eine revidierte, konsistente Analyse.
	
	\section{Grundlegende Konzepte der YUCT}
	
	\subsection{Koordinationseffizienz \(\Keff\)}
	
	In der YUCT wird jedes interagierende System durch ein Wörterbuch \(D\) (gemeinsame Protokolle, Zustände) und einen Index \(i\) charakterisiert, der eine bestimmte koordinierte Aktion aktiviert. Die Effizienz der Koordination ist
	\begin{equation}
		\Keff = \frac{H(\text{Aktionen})}{H(\text{Index})},
	\end{equation}
	wobei \(H\) die Shannon-Entropie bezeichnet. Für Quantensysteme kann \(\Keff\) extrem groß werden und im Grenzfall perfekter Verschränkung gegen Unendlich gehen.
	
	\subsection{Universelles Skalierungsgesetz}
	
	Fluktuationen und Fehler gehorchen Gl.~\eqref{eq:scaling}. Dieses Gesetz ergibt sich aus der fraktalen Geometrie von Wörterbüchern und wurde empirisch in Systemen so unterschiedlich wie dem Neutronenzerfall und galaktischen Rotationskurven bestätigt \cite{Yakushev2026c}. Die logarithmische Form ist wesentlich: Sie spiegelt wider, dass die effektive Auflösung im Indexraum nur logarithmisch mit der Wörterbuchgröße wächst.
	
	\subsection{Temperaturabhängigkeit von \(\Keff\)}
	
	In vielen physikalischen Systemen nimmt die Koordinationseffizienz mit der Temperatur aufgrund thermischer Fluktuationen ab. Aus kritischen Phänomenen und der fraktalen Skalierung haben wir
	\begin{equation}
		\Keff(T) = K_0 \left(\frac{T}{T_0}\right)^{-\gamma}, \quad \gamma > 0,
		\label{eq:K_T}
	\end{equation}
	wobei \(\gamma\) ein materialabhängiger Exponent ist. Das Produkt \(\betac\gamma\) wurde unabhängig aus geophysikalischen Daten bestimmt: \(\betagamma = 0,20 \pm 0,03\) \cite{Yakushev2026b}. Dies impliziert \(\gamma = \betagamma / \betac = 0,20 / (2/3) = 0,30\). Dieser Wert ist konsistent mit allgemeinen Skalierungsargumenten für ungeordnete Systeme.
	
	\section{YUCT-Modell der Photon-Atom-Wechselwirkung}
	
	Betrachten wir einen resonanten Photonenpuls, der auf ein Ensemble kalter Atome trifft. Die Atome wirken als \textbf{Wörterbuch} \(D\): Sie werden in einem bestimmten Quantenzustand präpariert (z.B. durch Laserkühlung und Fallen). Das Photon trägt einen \textbf{Index} – seine Wellenpaketform und Frequenz. Die Wechselwirkung führt zu einem koordinierten Zustand, in dem das Photon transmittiert, absorbiert oder gestreut werden kann.
	
	Im Experiment von Angulo et al. \cite{Angulo2024} ist die interessierende Größe die \textbf{atomare Anregungszeit} \(\tau_T\) für ein transmittiertes Photon. Aus der YUCT ist diese Zeit proportional zur Gruppenverzögerung \(\tau_g\). Die Weak-Value-Analyse von Thompson et al. \cite{Thompson2023} zeigt, dass \(\tau_T = \tau_g\) für ein transmittiertes Photon im relevanten Regime gilt.
	
	Der entscheidende Schritt ist die Verknüpfung von \(\tau_g\) mit \(\Keff\). Die Gruppenverzögerung entsteht aus der kohärenten Wechselwirkung zwischen dem Photon und dem atomaren Ensemble. Die Stärke dieser Wechselwirkung ist nicht proportional zu \(\Keff\) selbst (das enorm sein kann), sondern zur \textbf{Information}, die das Photon über den Atomzustand trägt. In der YUCT ist das relevante Informationsmaß der Logarithmus der Koordinationseffizienz, weil die Indexlänge \(\ell = \log_2 M\) die Anzahl der unterscheidbaren Zustände bestimmt. Daher postulieren wir:
	\begin{equation}
		|\tau_g| = \alpha_0 \, \ln\Keff,
		\label{eq:tau_log}
	\end{equation}
	wobei \(\alpha_0\) eine fundamentale Zeitkonstante mit der Dimension Zeit ist. Für die Rubidium-D2-Linie beträgt die natürliche Linienbreite \(\Gamma = 2\pi \times 6\) MHz, was eine charakteristische Zeit \(\tau_{\text{nat}} = 1/\Gamma \approx 26\) ns ergibt. Die Anpassung an den experimentellen Wert \(|\tau_g| \approx 30\) ns bei \(K_0 \sim 10^4\) liefert \(\alpha_0 = 30\,\text{ns} / \ln(10^4) = 30 / 9,21 \approx 3,26\) ns. Dies ist plausibel: Die elementare Wechselwirkungszeit liegt in der Größenordnung einiger Nanosekunden.
	
	\subsection{Temperaturabhängigkeit der Gruppenverzögerung}
	
	Einsetzen von Gl.~\eqref{eq:K_T} in Gl.~\eqref{eq:tau_log} ergibt
	\begin{equation}
		|\tau_g(T)| = \alpha_0 \ln\!\left(K_0 \left(\frac{T}{T_0}\right)^{-\gamma}\right)
		= \alpha_0 \ln K_0 - \alpha_0 \gamma \ln\!\left(\frac{T}{T_0}\right).
	\end{equation}
	Mit \(|\tau_g(T_0)| = \alpha_0 \ln K_0\) erhalten wir
	\begin{equation}
		|\tau_g(T)| = |\tau_g(T_0)| - \alpha_0 \gamma \ln\!\left(\frac{T}{T_0}\right).
		\label{eq:tau_T_log}
	\end{equation}
	Dies ist das zentrale Ergebnis: Der Betrag der Gruppenverzögerung nimmt \textbf{logarithmisch} mit steigender Temperatur ab. Für ein Temperaturverhältnis \(T/T_0 = 2\) beträgt die Änderung
	\begin{equation}
		\Delta |\tau_g| = \alpha_0 \gamma \ln 2.
	\end{equation}
	Mit \(\alpha_0 \approx 3,26\) ns und \(\gamma = 0,30\) finden wir \(\Delta |\tau_g| \approx 3,26 \times 0,30 \times 0,693 \approx 0,68\) ns. Die relative Änderung bezogen auf den Basiswert \(|\tau_g(T_0)| \approx 30\) ns beträgt \(0,68 / 30 \approx 2,3\%\). Dies liegt immer noch innerhalb der statistischen Unsicherheit des Experiments (\(\pm 0,22\) in Einheiten von \(\tau_0\), d.h. etwa \(\pm 5,7\) ns). Der Effekt liegt jedoch an der Nachweisgrenze. Noch wichtiger ist, dass die Änderung \textbf{logarithmisch} ist, nicht eine Potenzfunktion, und ihre Größe ist weit kleiner als die zuvor behaupteten 13\%.
	
	\subsection{Vergleich mit dem früheren (fehlerhaften) linearen Modell}
	
	In früheren Arbeiten wurde eine lineare Beziehung \(|\tau_g| \propto \Keff\) angenommen. Dies führte zu \(|\tau_g| \propto T^{-\gamma}\) und, mit \(\gamma = 0,30\), zu einer 13\%-Änderung bei einer Temperaturverdopplung. Diese Vorhersage beruhte auf einem Missverständnis der Rolle von \(\Keff\) in Quantensystemen. Die korrekte logarithmische Beziehung ergibt sich, weil \(\Keff\) ein Verhältnis von Entropien ist, keine direkte physikalische Feldamplitude. Das lineare Modell ist mit den informationstheoretischen Grundlagen der YUCT unvereinbar und wird hiermit zurückgenommen.
	
	\section{Numerische Konsistenzprüfung anhand vorhandener Daten}
	
	Das Experiment von Angulo et al. \cite{Angulo2024} liefert einen spezifischen Datenpunkt, der es uns ermöglicht, \(\alpha_0\) und \(K_0\) abzuschätzen. Für die 10-ns-Puls-Konfiguration mit optischer Tiefe OD~4 wurde eine atomare Anregungszeit von \(\tau_T = -1,14 \tau_0\) berichtet. Mit \(\tau_0 \approx 26\) ns ergibt sich \(|\tau_g| \approx 30\) ns. Unter der Annahme, dass \(\Keff(T_0) = K_0\) in der Größenordnung von \(10^4\) liegt (typisch für ein kaltes atomares Ensemble mit \(10^6\) Atomen und hoher Kohärenz), erhalten wir \(\alpha_0 = 30\,\text{ns} / \ln(10^4) \approx 3,26\) ns. Dieser Wert ist als charakteristische Zeitskala für die elementare Photon-Atom-Wechselwirkung plausibel. Es sind keine weiteren freien Parameter erforderlich.
	
	\section{Vergleich mit der Standard-Quantenoptik}
	
	Im konventionellen Rahmen der Quantenoptik ist die Gruppenverzögerung \(\tau_g\) für einen resonanten Puls durch die Ableitung der Transmissionsphase nach der Frequenz gegeben, \(\tau_g = d\phi/d\omega\). Für eine Lorentz-förmige atomare Antwort ergibt sich eine charakteristische Form, die aufgrund anormaler Dispersion auf dem Flügel der Resonanz negativ werden kann \cite{Milonni2004}. Die Standardtheorie sagt jedoch keine Abhängigkeit von \(\tau_g\) von der Temperatur der Atomwolke voraus, abgesehen von trivialen Effekten wie der Doppler-Verbreiterung (die linear in der Temperatur ist, nicht logarithmisch). Die beobachtete negative Zeitverzögerung wird als kohärenter Interferenzeffekt verstanden, der keinen Energieaustausch beinhaltet.
	
	YUCT geht über die Standardtheorie hinaus, indem sie die Größe der Verzögerung mit der Koordinationseffizienz \(\Keff\) verknüpft. Die logarithmische Abhängigkeit ist eine charakteristische Signatur: Wenn man \(\Keff\) durch Änderung der Atomdichte, der optischen Tiefe oder der Pulsdauer variieren könnte, sollte die Verzögerung wie \(\ln\Keff\) skalieren. Dies ist eine testbare Vorhersage, die nicht auf Temperaturvariationen angewiesen ist.
	
	\section{Überarbeitete experimentelle Vorschläge}
	
	Da die Temperaturabhängigkeit extrem schwach ist (logarithmisch, mit einer relativen Änderung von nur wenigen Prozent bei einer Temperaturverdopplung), ist es unwahrscheinlich, dass sie mit der aktuellen Präzision nachgewiesen werden kann. Wir schlagen daher alternative Experimente vor, die \(\Keff\) auf kontrolliertere Weise direkt variieren.
	
	\subsection{Variation der optischen Tiefe}
	
	Die Koordinationseffizienz \(\Keff\) ist proportional zur Anzahl der Atome im Ensemble und zur Wechselwirkungsstärke. Durch Verringerung der optischen Tiefe (z.B. durch Senken der Atomdichte) nimmt \(\Keff\) ab. Der Betrag der Gruppenverzögerung sollte dann \(|\tau_g| = \alpha_0 \ln\Keff\) folgen. Für eine Reduktion von \(\Keff\) um den Faktor 10 (z.B. von \(10^4\) auf \(10^3\)) würde die Änderung von \(|\tau_g|\) \(\alpha_0 (\ln 10^4 - \ln 10^3) = \alpha_0 \ln 10 \approx 3,26 \times 2,3026 \approx 7,5\) ns betragen, was etwa 25\% des Basiswerts entspricht und leicht messbar ist. Ein solches Experiment ist einfach: Wiederholen Sie die Messung von \(\tau_T\) für verschiedene Atomdichten (unter Beibehaltung aller anderen Parameter) und überprüfen Sie die logarithmische Skalierung.
	
	\subsection{Variation der Pulsdauer}
	
	Der Informationsgehalt des Photonenpulses (die Indexlänge) hängt von seiner zeitlichen Dauer und Bandbreite ab. Kürzere Pulse tragen mehr spektrale Information und haben daher ein größeres \(\Keff\). Durch systematische Variation der Pulsdauer von 1 ns bis 20 ns kann beobachtet werden, wie die Gruppenverzögerung mit \(\ln(\text{Bandbreite})\) skaliert. Dies bietet einen weiteren direkten Test des logarithmischen Gesetzes.
	
	\subsection{Temperaturabhängigkeit als Nulltest}
	
	Obwohl der vorhergesagte Temperatureffekt winzig ist, ist er nicht exakt Null. Eine Nullmessung (keine nachweisbare Änderung innerhalb der experimentellen Fehler) wäre mit unserem logarithmischen Modell konsistent, während das frühere lineare Modell eine klare 13\%-Änderung vorhergesagt hätte. Die vorhandenen Daten sprechen somit bereits gegen das lineare Modell, und zukünftige hochpräzise Experimente könnten schließlich die logarithmische Verschiebung nachweisen. Wir ermutigen Experimentatoren, beide zu testen.
	
	\section{Verbindung zur Photonenmassen-Vorhersage: Eine Konsistenzprüfung}
	\label{sec:photon_mass_consistency}
	
	Der YUCT-Rahmen hat auch eine blinde Vorhersage für die Ruhemasse des Photons hervorgebracht, die aus derselben Massenleiter \(m = m_e q^{N_f}\) mit \(q=(3/2)^{1/3}\) abgeleitet wurde. Der optimale Knoten, der mit experimentellen Obergrenzen konsistent ist, ist \(N_\gamma = -410\), was zu
	\[
	m_\gamma \approx 7,3 \times 10^{-19}~\text{eV} \quad (\text{siehe Anhang Photon Mass})
	\]
	führt.
	
	Man könnte sich fragen, ob eine solche winzige, von Null verschiedene Photonenmasse die Gruppenverzögerungsmessung in kalten Atomen beeinflussen könnte. Die Compton-Wellenlänge entsprechend \(m_\gamma\) ist
	\[
	\lambda_c = \frac{\hbar}{m_\gamma c} \approx 2,7\times10^{11}~\text{m},
	\]
	was um viele Größenordnungen größer ist als die Atomwolke (typische Größe \(\sim 1\) mm). Daher sind Ausbreitungseffekte aufgrund einer endlichen Photonenmasse in diesem Experiment völlig vernachlässigbar. Die negative Zeitverzögerung ist ein reiner Koordinationseffekt, unabhängig von der absoluten Massenskala.
	
	Umgekehrt stellt die Tatsache, dass ein einziger Rahmen korrekt sowohl
	\begin{itemize}
		\item die Existenz negativer Gruppenverzögerungen mit einer durch \(\ln\Keff\) bestimmten Größe (dieser Anhang),
		\item die Obergrenze für die Photonenmasse (\(m_\gamma < 10^{-18}\) eV) und ihren diskreten Knoten \(N_\gamma = -410\) (Anhang Photon Mass) vorhersagt,
	\end{itemize}
	eine leistungsfähige Konsistenzprüfung dar. Die beiden Vorhersagen betreffen stark unterschiedliche physikalische Bereiche (Quantenoptik vs. fundamentale Teilcheneigenschaften) und teilen dennoch dieselben zugrundeliegenden Konstanten \(q\), \(S_{\mathrm{odd}}/S_{\mathrm{even}}\) und das Konzept der Koordinationseffizienz \(\Keff\). Es sind keine ad-hoc-Anpassungen erforderlich, um sie kompatibel zu machen.
	
	Somit widerspricht das Experiment zur negativen Photonenverzögerung nicht der YUCT-Photonenmassen-Vorhersage; vielmehr dienen beide als unabhängige Pfeiler, die das vereinheitlichte Koordinationsparadigma stützen.
	
	\section{Rücknahme der früheren blinden Vorhersage}
	
	In einer früheren Version dieses Anhangs (V.2.0) sagten wir fälschlicherweise voraus, dass die Gruppenverzögerung wie \(T^{-0,20}\) skalieren würde, was zu einer 13\%-Änderung bei einer Temperaturverdopplung führen würde. Diese Vorhersage basierte auf einer falschen linearen Beziehung \(|\tau_g| \propto \Keff\). Die hier abgeleitete korrekte logarithmische Beziehung zeigt, dass der Temperatureffekt etwa zwei Größenordnungen kleiner ist. Daher ist die Behauptung eines 13\%-Effekts ungültig und wird formell zurückgenommen. Wir entschuldigen uns für das Versehen und danken den Mitautoren des mathematischen Patches für die Identifizierung der Inkonsistenz.
	
	Dennoch bleibt das YUCT-Kernkonzept – dass negative Gruppenverzögerungen eine Manifestation der Koordinationseffizienz sind – gültig. Die experimentelle Bestätigung negativer atomarer Anregungszeiten unterstützt bereits die Idee, dass Photon und Atom ein koordiniertes Wörterbuch teilen. Das korrigierte Skalierungsgesetz liefert nun einen verfeinerten theoretischen Rahmen für zukünftige Tests.
	
	\section{Fazit}
	
	Wir haben die YUCT-Analyse des Experiments zur negativen Photonenverzögerung korrigiert. Das universelle Fehlergesetz der YUCT ist logarithmisch, und dementsprechend muss die Beziehung zwischen Gruppenverzögerung und Koordinationseffizienz logarithmisch sein: \(|\tau_g| = \alpha_0 \ln\Keff\). Dies führt zu einer sehr schwachen Temperaturabhängigkeit (logarithmisch, mit einer relativen Änderung von nur wenigen Prozent bei einer Temperaturverdopplung) und eliminiert den zuvor behaupteten 13\%-Effekt. Die fehlerhafte Vorhersage wird zurückgenommen. Alternative Tests, wie die Variation der optischen Tiefe oder der Pulsdauer, bieten viel empfindlichere Sonden für die logarithmische Skalierung. Die Existenz negativer Zeitverzögerungen bleibt eine bemerkenswerte Bestätigung der koordinierten Natur von Quantenwechselwirkungen, und YUCT bietet eine konsistente informationstheoretische Interpretation.
	
	\section*{Danksagung}
	
	Der Autor dankt den Mitautoren des mathematischen Patches für den Hinweis auf die Inkonsistenz zwischen dem linearen Ansatz und der logarithmischen Struktur der YUCT. Dank gilt auch den Teilnehmern des YUCT-Seminars für anregende Diskussionen und der Gruppe in Toronto (Aephraim Steinberg, Daniela Angulo und Kollegen) für ihr schönes Experiment.
	
	\section*{Datenverfügbarkeit}
	
	Alle Berechnungen und zusätzlichen Materialien sind verfügbar unter \url{https://github.com/Alexey-Yakushev-YUCT/YPSDC}.
	
	\begin{thebibliography}{99}
		
		\bibitem{Angulo2024}
		D. Angulo, K. Thompson, A. Jiao, H. Wiseman, A. Steinberg, ``Experimental evidence that a photon can spend a negative amount of time in an atom cloud,'' arXiv:2409.03680 [quant-ph] (2024).
		
		\bibitem{Nixon2024}
		V.-M. Nixon et al., ``Measuring the time transmitted photons spend as atomic excitations,'' American Physical Society DAMOP Conference, Session Y07 (June 2024).
		
		\bibitem{Yakushev2026a}
		A. V. Yakushev, ``Yakushev Unified Coordination Theory (YUCT),'' Zenodo, \url{https://doi.org/10.5281/zenodo.18444599} (2026).
		
		\bibitem{Yakushev2026b}
		A. V. Yakushev, ``Appendix P: Temperature Dependence of Gravity,'' Zenodo (2026).
		
		\bibitem{Yakushev2026c}
		A. V. Yakushev, ``Appendix L: Fractal Error Scaling,'' Zenodo (2026).
		
		\bibitem{Yakushev2026d}
		A. V. Yakushev, ``Appendix X: Generalization of Shannon Information Theory,'' Zenodo (2026).
		
		\bibitem{Thompson2023}
		K. Thompson et al., ``How much time does a photon spend as an atomic excitation before being transmitted?'' arXiv:2310.00432 (2023).
		
		\bibitem{Brillouin1960}
		L. Brillouin, \textit{Wave Propagation and Group Velocity} (Academic Press, 1960).
		
		\bibitem{Milonni2004}
		P. Milonni, \textit{Fast Light, Slow Light and Left-Handed Light} (CRC Press, 2004).
		
		\bibitem{Wiseman2023}
		H. M. Wiseman, A. M. Steinberg, and M. Hallaji, ``Obtaining a single-photon weak value from experiments using a strong coherent state,'' AVS Quantum Sci. \textbf{5}, 024401 (2023).
		
	\end{thebibliography}
	
\end{document}